\documentclass[11pt]{article}
\usepackage[margin=1in]{geometry}
\usepackage{amsmath,amssymb,amsthm}
\usepackage[colorlinks=true,urlcolor=blue,linkcolor=blue]{hyperref}
\newtheorem{theorem}{Theorem}
\newtheorem{lemma}[theorem]{Lemma}
\newcommand{\R}{\mathbb R}
\newcommand{\diag}{\operatorname{diag}}
\newcommand{\rank}{\operatorname{rank}}
\newcommand{\norm}[1]{\left\lVert#1\right\rVert}
\newcommand{\tr}{\operatorname{tr}}
\title{A single doubled edge on a degree-three unweighted graph\protect\\can require relative rank of order $1/\varepsilon$}
\author{Independent research extension: bounded-degree construction}
\date{15 November 2026}
\begin{document}
\maketitle

For a connected graph Laplacian $L$, a single-edge doubling $L'=L+bb^T$,
and $0<\alpha<1$, put
\[
 C_\alpha=(L^\dagger)^{\alpha/2}(L'^\alpha-L^\alpha)
                    (L^\dagger)^{\alpha/2}.
\]
Negative powers are zero on the common constant kernel. The source graph
has conductance one and vertex mass one everywhere.

\begin{theorem}\label{thm:main}
For each fixed $0<\alpha<1$, there are constants $c_\alpha>0$ and
$\varepsilon_\alpha>0$ such that, for every
$0<\varepsilon<\varepsilon_\alpha$, some finite connected simple
unweighted graph of maximum degree three and one of its edges satisfy
\[
 \#\{j:\lambda_j(C_\alpha)>\varepsilon\}
                 \ge \frac{c_\alpha}{\varepsilon}
\]
when that edge is doubled. Consequently the supremal strict count on
this degree-three class is $\Theta_\alpha(1/\varepsilon)$.
\end{theorem}

The theorem disproves a size-independent logarithmic-squared relative
rank bound for bounded-degree unweighted graphs. Its graph size is
unrestricted and can be very large. It does not assert that every
determinant-admissible relative spectrum is attained in this class.

\section*{The algebraic separated-mass input}
We use the following finite-dimensional lemma, proved independently in
the companion note \emph{Separated masses and reverse Schur pivots}.
It is recalled explicitly to distinguish it from the graph limit below.

\begin{lemma}[Separated masses and reverse pivots]\label{lem:pivots}
Let $K\succ0$, $K'=K+vv^T$ in dimension $m$, and let
\[
 D_\delta=\diag(\delta^{(m-1)/2},\ldots,\delta^{1/2},1),\qquad
 A_\delta=D_\delta KD_\delta,\quad
 B_\delta=D_\delta K'D_\delta.
\]
For $S_i=\{i,\ldots,m\}$ and $S_{m+1}=\varnothing$, set
\[
 \rho_i=
 \frac{\det K'_{S_i,S_i}/\det K_{S_i,S_i}}
      {\det K'_{S_{i+1},S_{i+1}}/\det K_{S_{i+1},S_{i+1}}},
 \qquad \det K_{\varnothing,\varnothing}=1.
\]
As $\delta\downarrow0$, the relative fractional eigenvalues converge,
with ordering, to $\rho_i^\alpha-1$. The result also holds for convergent
families $K_\delta\to K\succ0$ and $v_\delta\to v$, and is therefore
uniform on compact subsets of these data.
\end{lemma}

The proof diagonalizes $A_\delta$: its eigenvalues satisfy
$\lambda_i\sim\kappa_i\delta^{m-i}$, where $\kappa_i$ is the reverse
LDL pivot of $K$. The leverage mass on scale $i$ converges to
$\tau_i-\tau_{i+1}$, where
$\tau_i=v_{S_i}^TK_{S_i,S_i}^{-1}v_{S_i}$. The fractional rank-one
resolvent formula then gives the diagonal limit
$[(1+\tau_i)/(1+\tau_{i+1})]^\alpha-1$, while its off-diagonal entries
vanish at rate controlled by $(\lambda_i/\lambda_j)^{(1-\alpha)/2}$.
The determinant lemma identifies this ratio with $\rho_i$.

\section*{Reservoirs with bounded degree and a free port}

Fix $\gamma_0>0$ and simple cubic graphs of arbitrarily large admissible
even sizes whose combinatorial Laplacian gaps are at least $\gamma_0$.
Such graphs of every sufficiently large even size follow from
Bordenave's Theorem~1 \cite{Bordenave}, a proof of Friedman's theorem
for uniform simple regular graphs: take $d=3$ and adjacency bound
strictly below three. For example $\gamma_0=(3-2\sqrt2)/2$ is valid.
Subdivide one edge and use the new
vertex $p$ as a port. Its internal degree is two. The resulting
reservoir $H$ still has a uniform gap
\[
 \gamma=\frac1{2+6/\gamma_0}>0.
\]
Indeed, writing $Q(x)$ for its energy and $\mu$ for the mean over its
old vertices, the old graph energy is at most $2Q(x)$, and
\[
 \sum_{\rm old}(x_v-\mu)^2\le\frac{2Q(x)}{\gamma_0},\qquad
 (x_p-\mu)^2\le2Q(x)+\frac{4Q(x)}{\gamma_0}.
\]
The variance about the actual mean is no larger. Define the port
resistance to the reservoir mean by
\[
 r(H,p)=e_p^TL_H^\dagger e_p,\qquad 0\le r(H,p)\le R:=1/\gamma.
\]
Reservoir sizes approximate any prescribed positive mass ratios by
rounding to the nearest admissible even old size before subdivision.

If a reservoir mean is $V$ and a unit edge joins its port to a core
potential $x$, minimizing over mean-zero reservoir coordinates $z$ gives
\[
 \min_{z\perp\mathbf1}\bigl[z^TL_Hz+(x-V-e_p^Tz)^2\bigr]
             =\frac{(x-V)^2}{1+r(H,p)}.
\]
For an incident edge of conductance $w$ the denominator is
$w^{-1}+r(H,p)$. Thus static reservoir elimination is exactly a series
contact resistance, rather than an assumption of a perfectly
conducting reservoir.

\section*{A binary fanout with long terminal necks}

Let $m=2^k\ge2$ and fix an integer $K\ge4(R+1)$. Take a complete
binary tree with root $o$ and $m$ leaves. Each leaf is connected by a
simple path of length
\[
 \ell=Km
\]
to the port of its own reservoir $H_i$. Connect the root by one unit
edge $e$ to the port of an anchor reservoir $H_0$, and double exactly
$e$. All vertices have degree at most three, and $e$ is a bridge.

Let $r_i$ be the reservoir contact resistances. After eliminating
mean-zero reservoir coordinates, each terminal neck has resistance
$\ell+r_i$. When the anchor mean is fixed to zero, the root is
grounded by conductance
\[
 c_0=\frac1{1+r_0},\qquad
 c_0'=\frac1{1/2+r_0}
\]
before and after the update. Eliminate the finite core consisting of
the tree and neck interiors. Let $K(r),K'(r)$ be the resulting SPD
stiffness matrices on the $m$ terminal means. Their difference is
positive rank one by the Schur-complement rank-one identity. They vary
continuously over the compact cube $r\in[0,R]^{m+1}$ and remain SPD.

Order the terminals arbitrarily. Let $g_j$ be the root admittance,
excluding the anchor branch, when terminals $1,\ldots,j$ are grounded
through their necks and the other terminals are open. Then
\[
 g_0=0,\qquad 0\le g_j\le m/\ell=1/K.
\]

\begin{lemma}[Uniformly nontrivial increments]\label{lem:increments}
For every $1\le j\le m$,
\[
 g_j-g_{j-1}\ge\frac1{4Km}.
\]
\end{lemma}
\begin{proof}
Hold the root at potential one before attaching shunt $j$.
All potentials lie in $[0,1]$, so each existing grounded terminal
carries current at most $1/\ell$. An edge at tree depth $h$ on the
path to terminal $j$ carries at most
$(m/2^h)/\ell=1/(K2^h)$. Consequently that leaf has open-circuit
potential $v_j\ge1-1/K\ge3/4$.

Let $\eta_j$ be the effective resistance of this leaf to the root
when the root and existing shunts are grounded. The direct tree path
gives $\eta_j\le k$. Adding the new shunt and using the scalar
rank-one inverse identity gives exactly
\[
 g_j-g_{j-1}=\frac{v_j^2}{\ell+r_j+\eta_j}.
\]
Since $r_j\le R$, $k\le m$, and $K\ge4(R+1)$, the denominator is
at most $2Km$. The increment is therefore at least $9/(32Km)$,
which proves the asserted weaker bound.
\end{proof}

Define
\[
 f(g)=\frac{c_0'+g}{c_0+g},\qquad
 \beta=\frac{1}{2(R+1/2)(R+1)(2+1/K)(1+1/K)}>0.
\]
For $0\le g\le1/K$ and $0\le r_0\le R$,
\[
 -\frac{d}{dg}\log f(g)
 =\frac{c_0'-c_0}{(c_0'+g)(c_0+g)}\ge\beta.
\]
For $S_i=\{i,\ldots,m\}$, Schur complementation and the determinant
lemma at the root give
\[
 \frac{\det K'_{S_i,S_i}}{\det K_{S_i,S_i}}
                  =\frac{f(g_{i-1})}{f(g_m)}.
\]
To verify the identity, keep the core and terminals $S_i$ and ground
all other terminals. The full determinant ratio after increasing the
root ground from $c_0$ to $c_0'$ is $f(g_{i-1})$. The eliminated
core-only determinant ratio is $f(g_m)$. Their quotient is the
terminal Schur determinant ratio.

The reverse pivot ratios in Lemma~\ref{lem:pivots} thus obey
\[
 \rho_i=\frac{f(g_{i-1})}{f(g_i)},\qquad
 \log\rho_i\ge\frac{\beta}{4Km}=\frac hm,
 \qquad h:=\frac\beta{4K}>0.
\]
Give terminal $i$ mass $M_i=\delta^{-(m-i)}$.
Lemma~\ref{lem:pivots} shows that all $m$ limiting fractional relative
eigenvalues are at least
\[
 \rho_i^\alpha-1\ge\alpha\log\rho_i\ge\frac{\alpha h}{m}.
\]
Compactness in $r$ makes this mass-separation choice uniform over all
allowable port parameters. Give the anchor an additional mass tending
to infinity. The mass-normalized connected stiffness pair converges to
zero on the anchor and the grounded SPD pair on the terminals.
Its nonzero eigenvalues stay bounded away from zero for fixed terminal
masses, uniformly in $r$. Hence positive and pseudoinverse negative
fractional powers converge, and this anchor-grounding step is also
uniform on the contact-parameter cube.

\section*{Passing from masses to ordinary unweighted reservoirs}

\begin{lemma}[Low relative compression of large reservoirs]
\label{lem:reservoir}
Fix a finite core and positive target masses $M_0,\ldots,M_m$.
Choose the unit-mass reservoirs above with sizes
$n_i(s)/s\to M_i$. Along a subsequence on which their port resistances
converge to $r_i$, the compression of the actual relative fractional
correction to the baseline low spectral subspace converges to the
relative fractional correction of the mass-normalized static Schur
pair with masses $M_i$ and contact resistances $r_i$.
The common kernel is included as one zero direction.
\end{lemma}
\begin{proof}
Let $P_s$ be the isometry whose columns are the normalized reservoir
indicators, and put all remaining coordinates, including the entire
core, in the complementary block. Both graph Laplacians satisfy
\[
 P_s^TL_sP_s=O(s^{-1}),\qquad
 P_s^TL_s(I-P_sP_s^T)=O(s^{-1/2}).
\]
Write the complementary coordinates as $(x,z)$, where $x$ is on the
fixed core and $z$ is mean zero on each reservoir. Its energy $E(x,z)$
is at least $\gamma\norm z^2$. Minimizing over $z$ gives the fixed
core form grounded through contact conductances at least $1/(1+R)$;
its smallest eigenvalue is some $\mu>0$, uniformly in the contact
parameters. Thus also $E(x,z)\ge\mu\norm x^2$, and consequently
\[
 E(x,z)\ge\tfrac12\min(\gamma,\mu)
                         (\norm x^2+\norm z^2).
\]
The complementary block therefore has a uniform positive lower bound.
Doubling the edge preserves this bound. All full Laplacian norms are
bounded by eight, independently of reservoir sizes.

In block notation write
\[
 L_s=\begin{pmatrix}A_s&B_s^T\\B_s&F_s\end{pmatrix},\qquad
 A_s=O(s^{-1}),\quad B_s=O(s^{-1/2}),\quad F_s\succeq cI.
\]
Writing $N_s=\diag(n_i(s))$, its zero-frequency Schur complement is
exactly
\[
 A_s-B_s^TF_s^{-1}B_s
                   =N_s^{-1/2}\mathcal K(r(s))N_s^{-1/2},
\]
where $\mathcal K$ is the connected stiffness matrix on all reservoir
means, including the anchor. The finite core and reservoir
mean-zero coordinates have all been eliminated in this identity.
For $0\le\lambda<c/2$, the frequency correction obeys
\[
 \norm{B_s^T[(F_s-\lambda I)^{-1}-F_s^{-1}]B_s}
                                          =O(\lambda/s).
\]
Thus there are exactly $m+1$ low eigenvalues, one zero and the others
of order $s^{-1}$, while every other eigenvalue stays bounded below.
The low spectral isometries differ from $P_s$ by $O(s^{-1/2})$.
Multiplying the low eigenvalues by $s$ gives the limiting
mass-normalized Schur pair.

It remains to justify the fractional relative metric. Compressing a
bounded high spectral part to either low subspace contributes
$O(s^{-1})$. Multiplication by $s^\alpha$ makes that contribution
$o(1)$, using $\alpha<1$. The low part gives the limiting positive
fractional power. In the baseline low eigenbasis, the pseudoinverse
factors scale as $s^{\alpha/2}$ and converge after this scaling, zero
on the shared kernel. Multiplying proves convergence of the relative
compression. No convergence assertion about the high relative
response is needed.
\end{proof}

We now choose the terminal mass separation and the anchor mass so
that the ideal relative pair has $m$ eigenvalues at least, say,
$3\alpha h/(4m)$, uniformly in $r$. Choose reservoir sizes with these
fixed positive mass ratios and pass to a subsequence of converging
port parameters. Lemma~\ref{lem:reservoir} and the min--max principle
then give, for sufficiently large finite sizes, an actual degree-three
unit-mass graph with
\[
 \lambda_m(C_\alpha)\ge\frac{\alpha h}{2m}.
\]
These are sequential finite approximation choices; no infinitely
massive vertex or infinite graph is used in the final witness.

\begin{proof}[Proof of Theorem~\ref{thm:main}]
For small $\varepsilon$, choose a power of two satisfying
\[
 \frac{\alpha h}{8\varepsilon}\le m
                  <\frac{\alpha h}{4\varepsilon}.
\]
The preceding construction has at least $m$ eigenvalues greater than
$\varepsilon$, proving the lower bound with $c_\alpha=\alpha h/8$.

For the upper bound on any connected unweighted graph, the leverage
of doubling one existing edge is $\tau=b^TL^\dagger b\le1$.
Positivity and the determinant lemma give
\[
 C_\alpha\succeq0,\qquad
 \det_{\mathbf1^\perp}(I+C_\alpha)=(1+\tau)^\alpha\le2^\alpha.
\]
If $q$ correction eigenvalues exceed $\varepsilon$, then
$(1+\varepsilon)^q<2^\alpha$, and hence
\[
 q<\frac{\alpha\log2}{\log(1+\varepsilon)}
                        =O_\alpha(1/\varepsilon).
\]
This matches the degree-three lower order.
\end{proof}

\paragraph{Source and scope.}
The only external existence input is the uniform-gap bounded-degree
reservoir family. The circuit, separated-scale, and fractional
compression arguments above are deductions for the relative-update
question. The construction does not give efficient graph size or an
algorithm for computing a relative truncation. Its low response uses
only a fixed positive fraction of the available determinant budget;
additional high response eigenvalues are allowed.

\begin{thebibliography}{9}
\bibitem{Bordenave} C.~Bordenave. A new proof of Friedman's second
eigenvalue theorem and its extension to random lifts.
\textit{Ann. Sci. Ecole Norm. Sup.} 53 (2020), 1393--1439. Theorem~1.
\url{https://arxiv.org/abs/1502.04482}.
\end{thebibliography}
\end{document}
